\chapter{Surface Integrals}
\section{Parametric Surfaces}
\begin{definitionenv}\label{def: parametric surface}
    Suppose $T\subseteq \R^2$ has nonempty and connected interior and $\widetilde T \supseteq T$ is an open set of $\R^2$. Let $\bm r: \widetilde T\to \R^3$ be a continuous vector function. Then we call the image $S=\bm r(T)$ of the restricted $\bm r\big|_T$ a \textbf{parametric surface} with respect to $(\bm r, T)$ and $(\bm r,T)$ a \textbf{parametrization} of $S$. Furthermore, if $\bm r\big|_T$ is injective, then we call $S$ a \textbf{simple parametric surface}.

    We typically regard $S$ as lying in $xyz-$space, and $\widetilde T$ and $T$ as lying in the $uv-$plane. Formally, there exists functions $X,Y,Z: \widetilde T\to \R$ such that
    $$ x = X(u,v),\quad y = Y(u,v),\quad z = Z(u,v). $$
    Then the vector equation for $S$ is
    $$ \bm r(u,v) = X(u,v)\bm i + Y(u,v)\bm j + Z(u,v)\bm k,\quad (u,v)\in \widetilde T \supseteq T. $$
\end{definitionenv}
\begin{remarkenv}
    Here, we need restriction of $\bm r$ from $\widetilde T$ to $T$ to ensure that the partial derivatives of $\bm r$ is well-defined on the boundary of $T$. Note the parametrization of a surface is not unique.
\end{remarkenv}
\begin{exampleenv}
    Particularly, if a surface $S$ is given explicitly by an equation of the form $z = f(x,y)$, then we can parametrize $S$ by
    $$ X(u,v) = u,\quad Y(u,v) = v,\quad Z(u,v) = f(u,v). $$
    Conversely, suppose that we can solve the equations
    $$ x = X(u,v),\quad y = Y(u,v) $$
    for $u$ and $v$ in terms of $x$ and $y$, then we can express $z$ as a function of $x$ and $y$ by substituting the expressions for $u$ and $v$ into the equation
    $$ z = Z(u,v). $$
\end{exampleenv}
\begin{definitionenv}
    Suppose $S$ is a parametric surface with parametrization $(\bm r:\widetilde T \to \R^3, T)$ where
    $$ \bm r(u,v) = X(u,v)\bm i + Y(u,v)\bm j + Z(u,v)\bm k,\quad (u,v)\in \widetilde T \supseteq T, $$
    where $X,Y,Z$ are differentiable on $\widetilde T$. Then the \textbf{fundamental vector product} of $\bm r$ is defined by the cross product
    $$ \dfrac{\partial \bm r}{\partial u}\times \dfrac{\partial \bm r}{\partial v} =
    \left(\pdv{X}{u}, \pdv{Y}{u}, \pdv{Z}{u}\right)\times \left(\pdv{X}{v}, \pdv{Y}{v}, \pdv{Z}{v}\right) =
    \begin{vmatrix}
        \bm i & \bm j & \bm k\\[6pt]
        \dfrac{\partial X}{\partial u} & \dfrac{\partial Y}{\partial u} & \dfrac{\partial Z}{\partial u}\\[6pt]
        \dfrac{\partial X}{\partial v} & \dfrac{\partial Y}{\partial v} & \dfrac{\partial Z}{\partial v}
    \end{vmatrix}. $$
    If $(u_0,v_0) \in T$ is a point such that the partial derivatives $\displaystyle \pdv{\bm r}{u}$ and $\displaystyle \pdv{\bm r}{v}$ both exist, and are both continuous at $(u_0,v_0)$ and
    $$ \left( \dfrac{\partial \bm r}{\partial u}\times \dfrac{\partial \bm r}{\partial v} \right) \Bigg|_{(u_0,v_0)} \neq \bm 0, $$
    then we say $\bm r$ is \textbf{regular} at $(u_0,v_0)$. If $\bm{r}$ is not regular at $(u_0, v_0) \in \widetilde{T}$, then we say that $\bm{r}$ \textbf{is singular} at $(u_0, v_0)$. If $\bm{r}$ is regular at every $(u_0, v_0) \in \widetilde{T}$, then we say that $(\bm{r}, T)$ is a \textbf{smooth parametric representation}.

    If $(x, y, z) \in S$ is a point such that, for every $(u_0, v_0) \in T$ with $\bm{r}(u_0, v_0) = (x, y, z)$, the parametric representation $\bm{r}$ is regular at $(u_0, v_0)$, then we say that $(x, y, z)$ \textbf{is a regular point} of $S$ \textbf{(with respect to $(\bm{r}, T)$)}. If $(x, y, z) \in S$ is not a regular point of $S$ (\textbf{with respect to $(\bm{r}, T)$}), then we say that $(x, y, z)$ \textbf{is a singular point} of $S$ (\textbf{with respect to $(\bm{r}, T)$}). The parametric surface $S$ is called \textbf{smooth} (\textbf{with respect to $(\bm{r}, T)$}) if all of its points are regular points (\textbf{with respect to $(\bm{r}, T)$}).
\end{definitionenv}
\begin{remarkenv}
    Note a regular point for one parametrization of a surface may be a singular point for another parametrization of the same surface. Thus the regularity \textbf{does} depend on the choice of parametrization. Moreover, $\bm r$ may not be injective. Thus the definition of a regular point requires that all points in the preimage of $(x, y, z)$ are regular points of $\bm r$.
\end{remarkenv}
\begin{exampleenv}\label{example: regularity of a surface given explicitly}
    Follow the setup in Definition~\ref{def: parametric surface} and suppose $f:\widetilde T\to \R$ is a differentiable function and $S$ is the surface given explicitly by the equation $z = f(x,y)$ for $(x,y) \in T \subseteq \widetilde T$. Then we can parametrize $S$ by
    $$ \bm r(u,v) = u\bm i + v\bm j + f(u,v)\bm k. $$
    Then the fundamental vector product of $\bm r$ is
    $$ \dfrac{\partial \bm r}{\partial u}\times \dfrac{\partial \bm r}{\partial v} 
    = -\pdv{f}{x}\bm i - \pdv{f}{y}\bm j + \bm k$$
    for all $(u,v)\in T$. Thus all the regular points are those points $(x,y,z)$ on $S$ such that one of $\displaystyle \pdv{f}{x}(x,y)$ and $\displaystyle \pdv{f}{y}(x,y)$ does not exist.
\end{exampleenv}
\begin{exampleenv}\label{example: regularity depends on parametrization}
    Suppose the hemisphere
    $$ S = \{(x,y,z)\in \R^3: x^2 + y^2 + z^2 = a^2, z \ge 0\} $$
    is parametrized by
    $$ \bm r(\theta, \phi) = a(\sin\theta\cos\phi,\ \sin\theta\sin\phi,\ \cos\theta),\quad (\theta,\phi)\in \R^2, $$
    and $T = \{(\theta, \phi): 0\le \phi < 2\pi, 0\le \theta \le \dfrac{\pi}{2}\}$. Then the norm of the fundamental vector product of $\bm r$ is
    $$ \left\|\pdv{\bm r}{\theta}\times \pdv{\bm r}{\phi} \right\|= a^2 \left\| -\sin\theta
(\sin\theta\cos\phi,\ \sin\theta\sin\phi,\ \cos\theta)\right\|
    = a^2\sin\theta. $$
    Note that, with respect to this parametrization, the parameter values $(0,\phi)$ are singular. However, parametrizing $S$ by
    $$ \bm R(u,v) = \begin{cases}
        (u, v, \sqrt{a^2-u^2-v^2}),\quad &(u,v)\in T'\\
        (u, v, 0), \quad &(u,v)\in \R^2\setminus T' 
    \end{cases}, $$
    where $T' = \{(u,v): u^2 + v^2 \le a^2\}$, we have
    $$ \pdv{\bm R}{u} \left( 0,0 \right) \times \pdv{\bm R}{v} \left( 0,0 \right) =\left( 1,0,0 \right) \times \left( 0,1,0 \right) = \left( 0,0,1 \right) \ne \bm 0, $$
    showing that the north pole is a regular point of $S$. Thus the regularity of a point depends on the choice of parametrization. In this case, the set of singular points is $\{ (x,y,0) : x^2 + y^2 = a^2 \}$ since the partial derivatives of $\bm R$ fails to exist and be continuous at these points.(Note they are regular points with respect to the first parametrization.)
\end{exampleenv}
\begin{notationenv}
    Expanding the fundamental vector product by the first row, we have
    $$ \pdv{\bm r}{u}\times \pdv{\bm r}{v} = \left( \dfrac{\partial Y}{\partial u} \dfrac{\partial Z}{\partial v} - \dfrac{\partial Y}{\partial v} \dfrac{\partial Z}{\partial u} \right) \bm i - \left( \dfrac{\partial X}{\partial u} \dfrac{\partial Z}{\partial v} - \dfrac{\partial X}{\partial v} \dfrac{\partial Z}{\partial u} \right) \bm j + \left( \dfrac{\partial X}{\partial u} \dfrac{\partial Y}{\partial v} - \dfrac{\partial X}{\partial v} \dfrac{\partial Y}{\partial u} \right) \bm k. $$
    For convenience, we adopt the following notation
    $$ \Jac{Y}{Z}{u}{v} := \dfrac{\partial Y}{\partial u} \dfrac{\partial Z}{\partial v} - \dfrac{\partial Y}{\partial v} \dfrac{\partial Z}{\partial u} = \det \begin{pmatrix}
        \dfrac{\partial Y}{\partial u} & \dfrac{\partial Y}{\partial v}\\[6pt]
        \dfrac{\partial Z}{\partial u} & \dfrac{\partial Z}{\partial v} 
    \end{pmatrix}.$$
    Then we can write
    $$ \pdv{\bm r}{u}\times \pdv{\bm r}{v} = \Jac{Y}{Z}{u}{v} \bm i - \Jac{X}{Z}{u}{v} \bm j + \Jac{X}{Y}{u}{v} \bm k. $$
\end{notationenv}
\begin{propositionenv}\label{prop: fundamental vector product is normal to the surface}
    Let $S$ be a parametric surface with parametrization $(\bm r: \widetilde T \to \R^3, T)$ where
    $$\begin{array}{r r c l}
        \bm r: & \widetilde T & \longrightarrow & S \subseteq \R^3 \\
        & (u,v) & \longmapsto & \left( X(u,v), Y(u,v), Z(u,v) \right)
    \end{array}$$
    and $C\subseteq T \subseteq \widetilde T$ be a smooth curve with parametrization $$ \begin{array}{r r c l}
        \bm \alpha: & [a,b] & \longrightarrow & C \subseteq T \subseteq \widetilde T \\
        & t & \longmapsto & (u(t), v(t))
    \end{array}
    $$
    Then the composite $\bm \rho = \bm r\big|_C \circ \bm \alpha : [a,b]\to \bm r(C) \subseteq S$ is a smooth parametrization of a curve $\bm r(C)$ on $S$. Moreover, for each $t\in [a,b]$, the vectors
    $$ \dv{\bm \rho}{t},\quad \left( \pdv{\bm r}{u}\times \pdv{\bm r}{v} \right)\Bigg|_{(u,v)=\bm \alpha(t)}$$
    are orthogonal to each other.
\end{propositionenv}
\begin{proofenv}
    We adopt the notation in the statement of the proposition. Then the composite $\bm \rho$ is given by
    $$ \bm \rho(t) = X(u(t), v(t))\bm i + Y(u(t), v(t))\bm j + Z(u(t), v(t))\bm k. $$
    Since $\bm \alpha$ and $\bm r$ are both continuously differentiable, $\bm \rho$ is also continuously differentiable. Thus $\bm r(C)$ is a smooth curve. Differentiating $\bm \rho$ with respect to $t$, we have
    $$ \begin{aligned}
        \dv{\bm \rho}{t} = u'(t)\left( \pdv{X}{u} \bm i + \pdv{Y}{u} \bm j + \pdv{Z}{u} \bm k \right) + v'(t)\left( \pdv{X}{v} \bm i + \pdv{Y}{v} \bm j + \pdv{Z}{v} \bm k \right)
        = u'(t)\pdv{\bm r}{u} + v'(t)\pdv{\bm r}{v}.
    \end{aligned} $$
    Note the fundamental vector product $\displaystyle \pdv{\bm r}{u}\times \pdv{\bm r}{v}$ is normal to both $\displaystyle \pdv{\bm r}{u}$ and $\displaystyle \pdv{\bm r}{v}$. Thus $\displaystyle \pdv{\bm r}{u}\times \pdv{\bm r}{v}$ is normal to $\displaystyle \dv{\bm \rho}{t}$.
\end{proofenv}
\begin{remarkenv}
    By Proposition~\ref{prop: fundamental vector product is normal to the surface}, the fundamental vector product $\displaystyle \pdv{\bm r}{u}\times \pdv{\bm r}{v}$ is said to be normal to the surface $S = \bm r(T)$. By definition, at each regular point $P\in \bm r(T)$,
    $$ \left( \pdv{\bm r}{u}\times \pdv{\bm r}{v} \right) \Bigg|_P \ne \bm 0, $$
    and the plane passing through $P$ with normal vector $\displaystyle \left( \pdv{\bm r}{u}\times \pdv{\bm r}{v} \right)\Bigg|_P$ is called the \textbf{tangent plane} to $\bm r(T)$ at $P$ with respect to $(\bm r, T)$. 
\end{remarkenv}
\begin{definitionenv}
    Suppose $T$ is a region in the $uv$-plane and $S\subseteq \R^3$ be a parametric surface with parametrization $(\bm r: \widetilde T\to \R^3, T)$. The \textbf{area} of $S$ is defined by the double integral
    $$ a(S) = \iint_T \left\| \pdv{\bm r}{u}\times \pdv{\bm r}{v} \right\| \dd u \dd v. $$
    Moreover, if the boundary of $T$ is content zero, then we only require differentiability of $\bm r$ on $\operatorname{int} T$ and define $a(S)$ as the corresponding double integral over $\operatorname{int} T$.
\end{definitionenv}
\begin{remarkenv}
    It's possible that the area of a parametric surface $S=\bm r(T)$ may not exist. However, if $T$ is a bounded region and $\bm r$ is continuously differentiable, then $a(S)$ exists. Later we'll see that under some natural conditions, the area of a parametric surface is independent of the choice of parametrization $\bm r$.
\end{remarkenv}
\begin{exampleenv}
    Following the parametrization of the surface $S$ in Example~\ref{example: regularity of a surface given explicitly}, we have
    $$
    a(S) = \iint_T \sqrt{1 + \left( \pdv{f}{x} \right)^2 + \left( \pdv{f}{y} \right)^2} \dd x \dd y.
    $$
    In particular, when $S$ lies in a plane parallel to the $xy$-plane, i.e., $f(x,y) = c$ for some constant $c$, we have
    $$ a(S) = \iint_T 1 \dd x \dd y, $$
    which agrees with our definition of area for a region in the plane.
\end{exampleenv}
\begin{exampleenv}
    Following the parametrization of the semisphere $S$ in Example~\ref{example: regularity depends on parametrization}, we have
    $$ a(S) = \iint_T a^2\sin\theta \dd \theta \dd \phi = a^2 \int_0^{\pi/2} \sin\theta \dd \theta \int_0^{2\pi} \dd \phi = 2\pi a^2. $$
\end{exampleenv}
\begin{exampleenv}
    Let $a>b>0$ and consider the torus $S$ with vector equation
    $$ \bm r(u,v) = (a + b\cos v)\cos u\bm i + (a + b\cos v)\sin u\bm j + b\sin v\bm k,\quad (u,v)\in \R^2,$$
    where $T = \{(u,v): 0\le u < 2\pi, 0\le v < 2\pi\}$. Then the fundamental vector product of $\bm r$ is
    $$ \pdv{\bm r}{u}\times \pdv{\bm r}{v} = a(R+a\cos v)(\cos v\cos u,\cos v\sin u, \sin v\big).$$
    Thus the area of $S$ is given by
    $$ a(S) = \iint_T a(b + a\cos v) \dd u \dd v = b \int_0^{2\pi} \dd u \int_0^{2\pi} (a + b\cos v) \dd v = 4\pi^2 ab. $$
\end{exampleenv}

\section{Surface Integrals}
\begin{definitionenv}
    Let $S$ be a parametric surface with paramatrization $(\bm r: \widetilde T\to \R^3, T)$ where $\bm r$ is differentiable. Let $f: S\to \R$ be a bounded function. Then the \textbf{surface integral of $f$ over $S$} is defined by the double integral
    $$ \iint_S f(x,y,z) \dd S := \iint_T f(\bm r(u,v)) \left\| \pdv{\bm r}{u}\times \pdv{\bm r}{v} \right\| \dd u \dd v. $$
    Morevoer, if the boundary of $T$ is content zero, then we only require differentiability of $\bm r$ on $\operatorname{int} T$ and define the surface integral as the corresponding double integral over $\operatorname{int} T$.
\end{definitionenv}
\begin{notationenv}
    The surface integral of $f$ over $S$ is also denoted by
    $$ \iint_S f \dd S. $$
\end{notationenv}
\begin{remarkenv}
    To show that the definition of surface integral is well-defined, we need to verify that the integral does not depend on the choice of parametrization.
\end{remarkenv}
\begin{definitionenv}
    Let $(\bm r: \widetilde T\to \R^3, T)$ and $(\bm s: \widetilde U\to \R^3, U)$ be two parametrizations of the same surface $S = \bm r(T) = \bm s(U)$ such that $\bm r$ and $\bm s$ are both differentiable functions. If there exists a bijective and continuously differentiable function $\bm \phi: \widetilde U\to \widetilde T$ with nonvanishing Jacobian determinant such that $\bm s = \bm r\circ \bm \phi$ and $\bm \phi(U) = T$, then the parametrizations $(\bm r, T)$ and $(\bm s, U)$ are said to be \textbf{smoothly equivalent}.
\end{definitionenv}
\begin{theoremenv}
    Let $(\bm r: \widetilde T\to \R^3, T)$ and $(\bm s: \widetilde U\to \R^3, U)$ be two parametrizations of the same surface $S = \bm r(T) = \bm s(U)$ such that $\bm r$ and $\bm s$ are both differentiable functions. Let $f: S\to \R$ be a bounded function. Suppose $(\bm r, T)$ and $(\bm s, U)$ are smoothly equivalent. If one of the double integrals
    $$ \iint_T f(\bm r(u,v)) \left\| \pdv{\bm r}{u}\times \pdv{\bm r}{v} \right\| \dd u \dd v, \quad \text{or} \quad \iint_U f(\bm s(\xi,\eta)) \left\| \pdv{\bm s}{\xi}\times \pdv{\bm s}{\eta} \right\| \dd \xi \dd \eta $$
    exists, then both of them exist and they are equal.
\end{theoremenv}
\begin{proofenv}
    By definition, there exists a bijective and continuously differentiable function $\bm \phi: \widetilde U\to \widetilde T$ with nonvanishing Jacobian determinant such that $\bm s = \bm r\circ \bm \phi$ and $\bm \phi(U) = T$. Writing $\bm \phi(\xi, \eta) = (u(\xi,\eta), v(\xi,\eta))$, we have
    $$ \begin{cases}
        \dfrac{\partial \bm s}{\partial \xi} &= \dfrac{\partial \bm r}{\partial u}\dfrac{\partial u}{\partial \xi} + \dfrac{\partial \bm r}{\partial v}\dfrac{\partial v}{\partial \xi},\\[6pt]
        \dfrac{\partial \bm s}{\partial \eta} &= \dfrac{\partial \bm r}{\partial u}\dfrac{\partial u}{\partial \eta} + \dfrac{\partial \bm r}{\partial v}\dfrac{\partial v}{\partial \eta}.
    \end{cases} $$
    Thus the fundamental vector product of $\bm s$ is given by
    $$ \pdv{\bm s}{\xi}\times \pdv{\bm s}{\eta} = \left( \pdv{u}{\xi}\pdv{v}{\eta} - \pdv{u}{\eta}\pdv{v}{\xi} \right) \left( \pdv{\bm r}{u}\times \pdv{\bm r}{v} \right). $$
    Since the Jacobian determinant of $\bm \phi$ is nonvanishing, we have
    $$ \left\| \pdv{\bm s}{\xi}\times \pdv{\bm s}{\eta} \right\| = \left| \pdv{u}{\xi}\pdv{v}{\eta} - \pdv{u}{\eta}\pdv{v}{\xi} \right| \left\| \pdv{\bm r}{u}\times \pdv{\bm r}{v} \right\| = \left| \Jac{u}{v}{\xi}{\eta} \right| \left\| \pdv{\bm r}{u}\times \pdv{\bm r}{v} \right\|. $$
    Hence, by the change of variables formula for double integrals, if one of the following integrals exists, then both of them exist and they are equal:
    $$ \begin{aligned}
        &\iint_T f(\bm r(u,v)) \left\| \pdv{\bm r}{u}\times \pdv{\bm r}{v} \right\| \dd u \dd v\\
        =& \iint_U f(\bm r(u(\xi,\eta), v(\xi,\eta))) \left\| \pdv{\bm r}{u}\times \pdv{\bm r}{v} \right\| \left| \Jac{u}{v}{\xi}{\eta} \right| \dd \xi \dd \eta\\
        =& \iint_U f(\bm s(\xi,\eta)) \left\| \pdv{\bm s}{\xi}\times \pdv{\bm s}{\eta} \right\| \dd \xi \dd \eta.
    \end{aligned} $$
\end{proofenv}
\begin{remarkenv}
    In above discussions, we define the surface integral of a scalar function over a parametric surface. We can also define the surface integral of a vector field over a parametric surface.
\end{remarkenv}
\begin{definitionenv}
    Let $(\bm r: \widetilde T\to \R^3, T)$ be a simple and smooth parametrization of a parametric surface $S$. Then restircting $\bm r$ to $T$ gives a bijective function $\bm r\big|_T: T\to S$. Since $\bm r\big|_T$ is bijective, we define the \textbf{unit normal vector} of $S$ at $\bm r(u,v)$ by
    $$ \bm n(x,y,z) = \left.\frac{\dfrac{\partial \bm r}{\partial u}\times \dfrac{\partial \bm r}{\partial v}}{\left\| \dfrac{\partial \bm r}{\partial u}\times \dfrac{\partial \bm r}{\partial v} \right\|} \right|_{(u,v) = \bm r^{-1}(x,y,z)}$$
    having the same direction as the fundamental vector product.

    Let $\bm F: \widetilde S\to \R^3$ be a continuous vector function where $\widetilde S$ is an open subset of $\R^3$ containing $S$. Then the \textbf{surface integral of $\bm F$ over $S$}, also called the \textbf{flux of $\bm F$ through $S$}, is defined by the scalar surface integral
    $$ \iint_S \bm F \cdot \dd \bm S := \iint_S  \left(\bm F \cdot \bm n \right) \dd S.$$
\end{definitionenv}
\begin{remarkenv}
    We write the components of $\bm F$ and $\bm r$ as follows:
    $$ \begin{aligned}
        \bm F(x,y,z) &= P(x,y,z)\bm i + Q(x,y,z)\bm j + R(x,y,z)\bm k,\\
        \bm r(u,v) &= X(u,v)\bm i + Y(u,v)\bm j + Z(u,v)\bm k.
    \end{aligned}
    $$
    Then the surface integral of $\bm F$ over $S$ can be written as
    $$ \iint_S \bm F \cdot \dd \bm S = \iint_T \left( P\Jac{Y}{Z}{u}{v} + Q\Jac{Z}{X}{u}{v} + R\Jac{X}{Y}{u}{v} \right) \dd u \dd v. $$
    It is common to write
    $$
    \begin{aligned}
        \iint_T P(\bm r(u,v)) \Jac{Y}{Z}{u}{v} \dd u \dd v &=: \iint_S P \dd y \wedge \dd z,\\
        \iint_T Q(\bm r(u,v)) \Jac{Z}{X}{u}{v} \dd u \dd v &=: \iint_S Q \dd z \wedge \dd x,\\
        \iint_T R(\bm r(u,v)) \Jac{X}{Y}{u}{v} \dd u \dd v &=: \iint_S R \dd x \wedge \dd y.
    \end{aligned}
    $$
    Thus, we have
    $$ \iint_S \bm F \cdot \dd \bm S = \iint_S P \dd y \wedge \dd z + Q \dd z \wedge \dd x + R \dd x \wedge \dd y. $$
\end{remarkenv}
\begin{exampleenv}
    Following the parametrization of the surface $S$ in Example~\ref{example: regularity of a surface given explicitly}, we have
    $$ \pdv{\bm r}{u}\times \pdv{\bm r}{v} = -\pdv{f}{x}\bm i - \pdv{f}{y}\bm j + \bm k. $$
    Let $\widetilde S$ be an open set of $\R^3$ containing $S$, and let $\bm F: \widetilde S\to \R^3$ be a continuous vector field given by
    $$ \bm F(x,y,z) = P(x,y,z)\bm i + Q(x,y,z)\bm j + R(x,y,z)\bm k. $$
    Then the surface integral of $\bm F$ over $S$ is given by
    $$ \iint_S \bm F \cdot \dd \bm S = \iint_T \left( -P(x,y,f(x,y))\pdv{f}{x} - Q(x,y,f(x,y))\pdv{f}{y} + R(x,y,f(x,y)) \right) \dd x \dd y. $$
\end{exampleenv}

\section{Curl and Divergence of Vector Fields}
\begin{definitionenv}
    Let $\bm F$ be a differentiable vector field given by
    $$ \bm F(x,y,z) = P(x,y,z)\bm i + Q(x,y,z)\bm j + R(x,y,z)\bm k. $$
    Then the \textbf{curl} of $\bm F$ is defined by
    $$ \Curl \bm F = \left( \pdv{R}{y} - \pdv{Q}{z} \right)\bm i + \left( \pdv{P}{z} - \pdv{R}{x} \right)\bm j + \left( \pdv{Q}{x} - \pdv{P}{y} \right)\bm k. $$
    The \textbf{divergence} of $\bm F$ is defined by
    $$  \Div \bm F = \pdv{P}{x} + \pdv{Q}{y} + \pdv{R}{z}. $$
\end{definitionenv}
\begin{remarkenv}
    It is often convenient to introduce the differential operator $\nabla$ defined by
    $$ \nabla = \pdv{}{x}\bm i + \pdv{}{y}\bm j + \pdv{}{z}\bm k. $$
    Then we can write
    $$ \Curl \bm F = \curl \bm F, \quad  \Div \bm F = \div \bm F. $$
    Here the dot and cross products are interpreted formally using the standard rules of vector algebra, treating $\nabla$ as a vector of differential operators acting on the components of $\bm F$. In particular, we have
    $$
    \Curl \bm F =\det \begin{pmatrix}
        \bm i & \bm j & \bm k\\[6pt]
        \displaystyle \pdv{x} & \displaystyle \pdv{y} & \displaystyle \pdv{z}\\[6pt]
        P & Q & R
    \end{pmatrix}.
    $$
\end{remarkenv}
\begin{propositionenv}
    Let $\bm F$ be a differentiable vector field given by
    $$ \bm F(x,y,z) = P(x,y,z)\bm i + Q(x,y,z)\bm j + R(x,y,z)\bm k. $$
    Let
    $$ J_{\bm F} = \begin{pmatrix}
        \dfrac{\partial P}{\partial x} & \dfrac{\partial P}{\partial y} & \dfrac{\partial P}{\partial z}\\[6pt]
        \dfrac{\partial Q}{\partial x} & \dfrac{\partial Q}{\partial y} & \dfrac{\partial Q}{\partial z}\\[6pt]
        \dfrac{\partial R}{\partial x} & \dfrac{\partial R}{\partial y} & \dfrac{\partial R}{\partial z}
    \end{pmatrix} $$
    be the Jacobian matrix of $\bm F$. Then we have
    $$ \Div \bm F = \tr(J_{\bm F}).$$
    Moreover, if $J_{\bm F}$ is symmetric, then $\Curl \bm F = \bm 0$.
\end{propositionenv}
\begin{exampleenv}
    Let
    $$S = \{(x,y,z)\in \R^3: x^2 + y^2 \ne 0\}$$
    and $\bm F: S\to \R^3$ be a vector field given by
    $$ \bm F(x,y,z) =  -\dfrac{y}{x^2 + y^2}\bm i + \dfrac{x}{x^2 + y^2}\bm j. $$
    Note $S$ is an open subset of $\R^3$ and $\bm F$ is a differentiable vector field on $S$. Then the Jacobian matrix of $\bm F$ is given by
    $$ J_{\bm F} = \begin{pmatrix}
        \dfrac{2xy}{(x^2 + y^2)^2} & \dfrac{y^2 - x^2}{(x^2 + y^2)^2} & 0\\[6pt]
        \dfrac{y^2-x^2}{(x^2 + y^2)^2} & -\dfrac{2xy}{(x^2 + y^2)^2} & 0\\[6pt]
        0 & 0 & 0
    \end{pmatrix}. $$
    Thus we have $\Div \bm F = 0$ and $\Curl \bm F = \bm 0$. 
\end{exampleenv}
\begin{propositionenv}
    Let $\bm F$ and $\bm G$ be two twicely differentiable vector fields and $f$ be a twice continuously differentiable scalar field. Then we have
    \begin{enumerate}[label=(\arabic*)]
        \item for any $\alpha, \beta \in \R$,
        $$ \Curl(\alpha\bm F + \beta\bm G) = \alpha\Curl \bm F + \beta\Curl \bm G, \quad \Div(\alpha\bm F + \beta\bm G) = \alpha\Div \bm F + \beta\Div \bm G; $$
        \item $$ \Div(f\bm F) = f \Div \bm F + \bm F \cdot \nabla f; $$
        \item $$ \Curl(\nabla f) = \bm 0, \quad \Div(\curl \bm F) = 0. $$
        \item $$ \Div (\nabla f) = \pdv[2]{f}{x} + \pdv[2]{f}{y} + \pdv[2]{f}{z} =: \Delta f, $$
        where $\Delta$ is called the \textbf{Laplacian} operator.
    \end{enumerate}
\end{propositionenv}

\section{Stoke's Theorem}
\begin{theoremenv}[Stoke's Theorem]
    Let $T\subseteq \R^2$ where $\operatorname{int} T$ is the inner region of a piecewise smooth Jordan curve $\Gamma$. Suppose $S$ is a simple smooth parametric surface with parametrization $\bm r: T\to \R^3$ such that $\bm r$ is a twice continuously differentiable function on an open set $\widetilde{T} \supseteq T \cup \Gamma$ and $\Gamma$ is the boundary of $T$. Let $C= \bm r(\Gamma)$ be the boundary of $S$ and $\bm F: S\to \R^3$ be a continuously differentiable function on an open set $\widetilde{S} \supseteq \R^3$ with $\widetilde{S} \supseteq \bm r(T \cup \Gamma)$. Suppose $\Gamma$ is parametrized by $\bm \gamma: [a,b]\to \Gamma$, traversing the Jordan curve $\Gamma$ in the counterclockwise direction, and $C$ is parametrized by $\bm \alpha = \bm r\big|_\Gamma \circ \bm \gamma: [a,b] \to \R^3$. Then we have
    $$ \iint_S (\curl \bm F) \cdot \dd \bm S = \oint_C \bm F \cdot \dd \bm \alpha. $$
\end{theoremenv}
\begin{proofenv}
    We write the components of $\bm F$ as
    $$ \bm F(x,y,z) = P(x,y,z)\bm i + Q(x,y,z)\bm j + R(x,y,z)\bm k. $$
    Then the surface integral of $\curl \bm F$ over $S$ is given by
    $$ \iint_S \curl \bm F \cdot \dd \bm S = \iint_T \left( \pdv{R}{y} - \pdv{Q}{z} \right) \dd y \wedge \dd z + \left( \pdv{P}{z} - \pdv{R}{x} \right) \dd z \wedge \dd x + \left( \pdv{Q}{x} - \pdv{P}{y} \right) \dd x \wedge \dd y. $$
    Moreover, the line integral of $\bm F$ over $C$ is given by
    $$ \oint_C \bm F \cdot \dd \bm \alpha = \int_C P \dd x + Q \dd y + R \dd z. $$
    Therefore, it suffices to show that
    $$ \begin{cases}
        &\displaystyle \iint_S -\pdv{P}{y} \dd x \wedge \dd y + \pdv{P}{z} \dd z \wedge \dd x = \int_C P \dd x\\[8pt]
        &\displaystyle \iint_S -\pdv{Q}{z} \dd y \wedge \dd z + \pdv{Q}{x} \dd x \wedge \dd y = \int_C Q \dd y\\[8pt]
        &\displaystyle \iint_S -\pdv{R}{x} \dd z \wedge \dd x + \pdv{R}{y} \dd y \wedge \dd z = \int_C R \dd z
    \end{cases}
    $$
    and add them together. Note all of these integrals exist because
    \begin{itemize}
        \item $T \cup \Gamma \subseteq \widetilde{T}$, where $\bm r$ is a twice continuously differentiable function on $\widetilde{T}$ and $T \cup \Gamma$ is compact;
        \item $\bm r(T \cup \Gamma) \subseteq \widetilde{S}$, where $\bm F$ is a continuously differentiable function on $\widetilde{S}$ and $\bm r(T \cup \Gamma)$ is compact.
    \end{itemize}
    We only show the first equality holds since the other two are analogous. Note we can write the surface integral over $S$ as a double integral over $T$ and apply Green's theorem to write the double integral as a line integral over $\Gamma$. First, we have
    $$ \iint_S -\pdv{P}{y} \dd x \wedge \dd y + \pdv{P}{z} \dd z \wedge \dd x = \iint_T -\pdv{P}{y} \Jac{X}{Y}{u}{v} + \pdv{P}{z} \Jac{Z}{X}{u}{v} \dd u \dd v. $$
    Let
    $$ U = \bm r^{-1}(\widetilde S) := \{(u,v)\in \widetilde T: \bm r(u,v) \in \widetilde S\}. $$
    Since $\bm r: \widetilde T \to \R^3$ is continuous and $\widetilde S\supseteq \bm r(T \cup \Gamma)$, we know $U\subseteq \R^2$ is open and $U \supseteq T \cup \Gamma$. We restrict $\bm r$ to $U$:
    $$ \begin{array}{r r c l}
        \bm r\big|_U: & U & \longrightarrow & \widetilde S \subseteq \R^3\\
        & (u,v) & \longmapsto & (X(u,v), Y(u,v), Z(u,v))
    \end{array} $$
    and define
    $$ \begin{array}{r r c l}
        p:= P \circ \bm r\big|_U: & U & \longrightarrow & \R\\
        & (u,v) & \longmapsto & P(X(u,v), Y(u,v), Z(u,v))
    \end{array}
    $$
    which is a continuously differentiable function on $U$. We claim that
    $$ -\pdv{P}{y} \Jac{X}{Y}{u}{v} + \pdv{P}{z} \Jac{Z}{X}{u}{v} = \pdv{u} \left( p \pdv{X}{v} \right) - \pdv{v} \left( p \pdv{X}{u} \right).$$
    To see this, we first expand the right-hand side:
    $$ \pdv{u} \left( p \pdv{X}{v} \right) - \pdv{v} \left( p \pdv{X}{u} \right) = \pdv{p}{u} \pdv{X}{v} - \pdv{p}{v} \pdv{X}{u}.$$
    Then we apply the chain rule to write $\displaystyle \pdv{p}{u}$ and $\displaystyle \pdv{p}{v}$ as follows:
    $$ \begin{cases}
        \displaystyle \pdv{p}{u} = \pdv{P}{x}\pdv{X}{u} + \pdv{P}{y}\pdv{Y}{u} + \pdv{P}{z}\pdv{Z}{u},\\[6pt]
        \displaystyle \pdv{p}{v} = \pdv{P}{x}\pdv{X}{v} + \pdv{P}{y}\pdv{Y}{v} + \pdv{P}{z}\pdv{Z}{v}.
    \end{cases} $$
    Substituting these expressions into the previous expansion, we have
    $$ \begin{aligned}
        \pdv{u} \left( p \pdv{X}{v} \right) - \pdv{v} \left( p \pdv{X}{u} \right) &= \left( \pdv{P}{x}\pdv{X}{u} + \pdv{P}{y}\pdv{Y}{u} + \pdv{P}{z}\pdv{Z}{u} \right) \pdv{X}{v} - \left( \pdv{P}{x}\pdv{X}{v} + \pdv{P}{y}\pdv{Y}{v} + \pdv{P}{z}\pdv{Z}{v} \right) \pdv{X}{u}\\[6pt]
        &= \pdv{P}{y}\left( \pdv{Y}{u}\pdv{X}{v} - \pdv{Y}{v}\pdv{X}{u} \right) + \pdv{P}{z}\left( \pdv{Z}{u}\pdv{X}{v} - \pdv{Z}{v}\pdv{X}{u} \right)\\[6pt]
        &= -\pdv{P}{y} \Jac{X}{Y}{u}{v} + \pdv{P}{z} \Jac{Z}{X}{u}{v}.
    \end{aligned} $$
    Hence, we have
    $$ \iint_S -\pdv{P}{y} \dd x \wedge \dd y + \pdv{P}{z} \dd z \wedge \dd x = \iint_T \left( \pdv{u} \left( p \pdv{X}{v} \right) - \pdv{v} \left( p \pdv{X}{u} \right) \right) \dd u \dd v. $$
    Note $p$ is continuously differentiable on $U$ and $X$ is twice continuously differentiable on $U$. Therefore, $\displaystyle p \pdv{X}{v}$ and $\displaystyle p \pdv{X}{u}$ are continuously differentiable on $U\supseteq T \cup \Gamma$. Thus, by Green's theorem, we have
    $$ \iint_T \left( \pdv{u} \left( p \pdv{X}{v} \right) - \pdv{v} \left( p \pdv{X}{u} \right) \right) \dd u \dd v = \int_\Gamma p \pdv{X}{v} \dd v + p \pdv{X}{u} \dd u. $$
    Finally, we write $\bm \gamma = u \bm i + v \bm j$ and by definition of the line integral and chain rule, we have
    $$
    \int_\Gamma p \pdv{X}{v} \dd v + p \pdv{X}{u} \dd u = \int_a^b p(\bm \gamma (t)) \left( \left. \pdv{X}{u}\right|_{\bm \gamma(t)} u'(t) + \left. \pdv{X}{v}\right|_{\bm \gamma(t)} v'(t) \right) \dd t = \int_a^b p(\bm \gamma(t)) \dv{t} (X\circ \bm \gamma)(t) \dd t.
    $$
    Moreover, since $\bm \alpha = \bm r\big|_\Gamma \circ \bm \gamma$ on $[a,b]$ and $p = P \circ \bm r$ on $\Gamma$, we have
    $$ p \circ \bm \gamma = P \circ \bm r\big|_\Gamma \circ \bm \gamma = P \circ \bm \alpha, $$
    and thus
    $$
    \int_a^b p(\bm \gamma(t)) \dv{t} (X\circ \bm \gamma)(t) \dd t = \int_a^b P(\bm \alpha(t)) \dv{t} (X\circ \bm \gamma)(t) \dd t.
    $$
    Furthermore, $\bm \alpha : [a,b] \to C$ parametrizes $C$ and $X\circ \gamma$ is precisely the $x$-component of $\alpha$. Hence, we have
    $$ \int_a^b P(\bm \alpha(t)) \dv{t} (X\circ \bm \gamma)(t) \dd t = \int_C P \dd x. $$
    Therefore, we have shown that
    $$ \iint_S -\pdv{P}{y} \dd x \wedge \dd y + \pdv{P}{z} \dd z \wedge \dd x = \int_C P \dd x. $$
    Adding the other two equalities together, we finish the proof of Stoke's theorem.
\end{proofenv}
\begin{exampleenv}
    Let $C$ be the curve of intersection of the sphere $x^2 + y^2 + z^2 = a^2$ and the plane $x+y+z=0$. We want to evaluate the line integral
    $$ \int_C y\dd x + z\dd y + x\dd z. $$
    Let
    $$ \begin{array}{r r c l}
        \bm F:& \R^3 & \longrightarrow & \R^3\\
        & (x,y,z) & \longmapsto & y\bm i + z\bm j + x\bm k
    \end{array}
    $$
    We will find a simple smooth parametrization $(\bm r: \widetilde T\to \R^3, T)$ of a surface $S$ such that $\partial S = C$ and apply Stoke's theorem to evaluate the line integral. 

    Note $S$ should be a disk in the $x+y+z=0$ plane. Let $\widetilde T = \R^2$ and $T = \{(u,v)\in \R^2: u^2 + v^2 \le a^2\} \subseteq \widetilde{T}$. Note $T$ is closed with nonempty connected interior, and its boundary $\Gamma = \{(u,v)\in \R^2: u^2 + v^2 = a^2\}$ is a smooth Jordan curve.

    We define
    $$ \begin{array}{r r c l}
        \bm r:& \widetilde T & \longrightarrow & \R^3\\
        & (u,v) & \longmapsto & u\bm b_1 + v\bm b_2
    \end{array} $$
    where $\{ \bm b_1, \bm b_2\}$ is an orthonormal basis for the plane $x+y+z=0$. In particular, we can take
    $$ \bm b_1 = \frac{1}{\sqrt{6}}(1,1,-2), \quad \bm b_2 = \frac{1}{\sqrt{2}}(1,-1,0). $$
    Then we have $\bm r(T) = S$ since
    $$
    \| \bm r(u,v) \| ^2 = u^2 + v^2 \le a^2.
    $$
    Letting $(x,y,z) = \bm r(u,v)$ be the corresponding point on the plane $x+y+z=0$ for any $(u,v)\in \widetilde T$, we have
    $$ (x,y,z) \in S \Leftrightarrow x^2+y^2+z^2 \le a^2 \Leftrightarrow u^2 + v^2 \le a^2 \Leftrightarrow (u,v) \in T. $$
    Moreover, we have
    $$
    \pdv{\bm r}{u} \times \pdv{\bm r}{v} = -\frac{1}{\sqrt{3}}(1,1,1).
    $$
    Thus, $(\bm r: \widetilde T\to \R^3, T)$ is a simple smooth parametrization of $S$. We parametrize the boundary $\Gamma$ of $T$ by
    $$ \begin{array}{r r c l}
        \bm \gamma:& [0,2\pi] & \longrightarrow & \Gamma\\
        & t & \longmapsto & a\cos t \bm i + a\sin t \bm j.
    \end{array} $$
    Then we parametrize $C$ by $\bm \alpha = \bm r\big|_\Gamma \circ \bm \gamma: [0,2\pi] \to C$ and we have $r(\partial T)=C.$
    The curl of $\bm F$ is given by
    $$ \curl \bm F = \left( \pdv{y} x - \pdv{z} z \right)\bm i + \left( \pdv{z} y - \pdv{x} x \right)\bm j + \left( \pdv{x} z - \pdv{y} y \right)\bm k = -\bm i - \bm j - \bm k. $$
    Thus, by Stoke's theorem, we have
    $$
    \int_C \bm F \cdot \dd \bm \alpha = \iint_S (\curl \bm F) \cdot \dd \bm S = \int_T (\curl \bm F) \big|_{(x,y,z)=\bm r(u,v)} \cdot \left( \pdv{\bm r}{u} \times \pdv{\bm r}{v} \right) \dd u \dd v = \int_T \sqrt3 \dd u \dd v = \sqrt3 \pi a^2.
    $$
\end{exampleenv}

\section{Orientable and Closed Surfaces}
\begin{definitionenv}
    A subset $S\subseteq \R^3$ is said to be a \textbf{piecewise smooth surface} if it can be expressed as a finite union
    $$ S = \bigcup_{i=1}^n S_i $$
    where each $S_i$ admits a simple smooth parametrization $(\bm r_i, T_i)$ and the intersection of any two distinct $S_i$ and $S_j$ is either empty, a finite set of points, or a finite number of smooth curves; we call these intersections \textbf{common boundaries}.
\end{definitionenv}
\begin{remarkenv}
    There are two parametrizations of the northern hemisphere $S$ of radius $a>0$ in Example~\ref{example: regularity depends on parametrization}. However, neither of them is a simple smooth parametrization of $S$. We can combine them to make a simple smooth parametrization of $S$ by the following example.
\end{remarkenv}
\begin{exampleenv}\label{example: orientable surface}
     Suppose $S$ is the northern hemisphere of radius $a>0$. Let
     $$ \cos \theta_0 = \frac{b}{a}, \quad 0< b < a,\quad 0 < \theta_0 < \frac{\pi}{2}. $$
    We parametrize the top cap $S_1$ of $S$ by
    $$ \begin{array}{r r c l}
        \bm r_1: & \widetilde T_1 & \longrightarrow & \R^3\\
        & (x,y) & \longmapsto & (x,y,\sqrt{a^2 - x^2 - y^2})
    \end{array}, $$
    where $\widetilde T_1 = \{(x,y)\in \R^2: x^2 + y^2 < a^2\}$ and $T_1 = \{(x,y)\in \R^2: x^2 + y^2 \le b^2\} \subseteq \widetilde T_1$. We parametrize the bottom part $S_2$ by
    $$ \begin{array}{r r c l}
        \bm r_2: & \widetilde T_2 & \longrightarrow & \R^3\\
        & (\theta, \phi) & \longmapsto & (a\sin\theta\cos\phi, a\sin\theta\sin\phi, a\cos\theta)
    \end{array}, $$
    where $\widetilde T_2 = \R^2$ and $T_2 = \{(\theta, \phi)\in \R^2: 0 \le \theta \le \theta_0, 0 \le \phi < 2\pi\} \subseteq \widetilde T_2$. 
    Then $S=S_1\cup S_2$ is piecewise smooth with respect to $(\bm r_1, T_1)$ and $(\bm r_2, T_2)$.
\end{exampleenv}
\begin{definitionenv}
    Let $S=\displaystyle \bigcup_{i=1}^n S_i$ be a piecewise smooth surface. We say $S$ is \textbf{orientable} if it is possible to choose simple smooth parametrizations $(\bm r_i, T_i)$ of $S_i$ such that, for every distince pair of adjacent patches $S_i$ and $S_j$, sharing a common boundary curve $C_{ij} = S_i \cap S_j$, the induced direction on $C$ from $S_i$ is opposite to that induced from $S_j$.
\end{definitionenv}
\begin{exampleenv}
    Note the northern hemisphere is orientable in Example~\ref{example: orientable surface}. However, the Möbius strip is not orientable.
\end{exampleenv}
\begin{definitionenv}
    An orientable piecewise smooth surface $S$ is said to be \textbf{closed} if its entire outer boundary is empty, i.e., $\partial S = \varnothing$.
\end{definitionenv}

\section{Gauss's Theorem}
\begin{theoremenv}
    Let $V$ be a closed and bounded solid in $\R^3$ whose boundary $S=\partial V$ is a piecewise smooth closed surface. Let $\bm n$ be the outer unit normal vector of $S$ pointing away from the inner region $V$. Suppose $\bm F$ is a continuously differentiable vector field on an open set containing $V$. Then we have
    $$ \iint_S \bm F \cdot \dd \bm S = \iiint_V \Div \bm F \dd V. $$
\end{theoremenv}
\begin{proofenv}
    We only prove the case when $V$ is convex. We write the components of $\bm F$ and $\bm n$ as
    $$
    \begin{cases}
        \bm F(x,y,z)&= P(x,y,z)\bm i + Q(x,y,z)\bm j + R(x,y,z)\bm k,\\
        \bm n(x,y,z)&= \cos(\alpha(x,y,z))\bm i + \cos(\beta(x,y,z))\bm j + \cos(\gamma(x,y,z))\bm k.
    \end{cases}
    $$
    Then the desired equality can be written as
    $$ \iiint_V \left( \pdv{P}{x} + \pdv{Q}{y} + \pdv{R}{z} \right) \dd V = \iint_S \left( P\cos\alpha + Q\cos\beta + R\cos\gamma \right) \dd S. $$
    Therefore, it suffices to show that
    $$ \begin{cases}
        &\displaystyle \iiint_V \pdv{P}{x} \dd V = \iint_S P\cos\alpha \dd S\\[8pt]
        &\displaystyle \iiint_V \pdv{Q}{y} \dd V = \iint_S Q\cos\beta \dd S\\[8pt]
        &\displaystyle \iiint_V \pdv{R}{z} \dd V = \iint_S R\cos\gamma \dd S
    \end{cases}
    $$
    We only prove the third equality since the other two are analogous. Since $V$ is convex, we can write
    $$ V = \{(x,y,z)\in \R^3: (x,y)\in D, g_1(x,y) \le z \le g_2(x,y)\} $$
    where $D\subseteq \R^2$ is a closed, connected and bounded region and $g_1, g_2: D \to \R$ are continuous functions. Then we have
    $$ \iiint_V \pdv{R}{z} \dd V = \iint_D \left( \int_{g_1(x,y)}^{g_2(x,y)} \pdv{R}{z} \dd z \right) \dd x \dd y = \iint_D R(x,y,g_2(x,y)) - R(x,y,g_1(x,y)) \dd x \dd y. $$
    Then, we evaluate the surface integral on the right-hand side by deviding $S$ into three parts:
    \begin{enumerate}
        \item the bottom surface $S_1$, where $z=g_1(x,y)$;
        \item the top surface $S_2$, where $z=g_2(x,y)$;
        \item the side surface $S_3$.
    \end{enumerate}
    On the bottom surface $S_1$, we have
    $$
    \bm r_1(x,y) = x\bm i + y\bm j + g_1(x,y)\bm k, \quad (x,y)\in D.
    $$
    Then the fundamental vector product points in the opposite direction of the outer unit normal vector $\bm n$ and we have
    $$
    \iint_{S_1} R\cos\gamma \dd S = -\iint_{S_1} R \dd x \wedge \dd y = -\iint_D R(x,y,g_1(x,y)) \dd x \dd y.
    $$
    On the top surface $S_2$, we have
    $$
    \bm r_2(x,y) = x\bm i + y\bm j + g_2(x,y)\bm k, \quad (x,y)\in D.
    $$
    Then the fundamental vector product points in the same direction of the outer unit normal vector $\bm n$ and we have
    $$
    \iint_{S_2} R\cos\gamma \dd S = \iint_{S_2} R \dd x \wedge \dd y = \iint_D R(x,y,g_2(x,y)) \dd x \dd y.
    $$
    On $S_3$, the fundamental vector product is perpendicular to $\bm k$ so $\gamma = \pi/2$ and we have
    $$
    \iint_{S_3} R\cos\gamma \dd S = 0.
    $$
    Summing up the three parts, we have
    $$ \iint_S R\cos\gamma \dd S = \iint_D R(x,y,g_2(x,y)) - R(x,y,g_1(x,y)) \dd x \dd y. $$
    This completes the proof of Gauss's theorem for convex solids.
\end{proofenv}
\begin{exampleenv}
    Let $V = [0,1]^3$ be the unit cube in $\R^3$ and $S = \partial V$ be the boundary of $V$. Let
    $$ \begin{array}{r r c l}
        \bm F:& \R^3 & \longrightarrow & \R^3\\
        & (x,y,z) & \longmapsto & x^2\bm i + y^2\bm j + z^2\bm k
    \end{array}
    $$
    Note the surface integral of $\bm F$ over $S$, $\displaystyle \iint_S \bm F \cdot \dd \bm S$, is the sum of the surface integrals over the six faces of the cube. It is quite annoying if we compute the surface integral directly over six faces. However, we can apply Gauss's theorem to compute the surface integral by evaluating the triple integral of $\Div \bm F$ over $V$. We have
    $$ \Div \bm F = \pdv{x} x^2 + \pdv{y} y^2 + \pdv{z} z^2 = 2x + 2y + 2z. $$
    Thus, by Gauss's theorem, we have
    $$ \iint_S \bm F \cdot \dd \bm S = \iiint_V \Div \bm F \dd V = \iiint_{[0,1]^3} 2x + 2y + 2z \dd x \dd y \dd z = 3. $$
\end{exampleenv}
\begin{corollaryenv}
    Let $V(t)$ be the solid sphere of radius $t>0$ centered at a point $\bm a\in \R^3$ and $S(t) = \partial V(t)$ be the boundary of $V(t)$. Suppose $\bm F$ is a continuously differentiable vector field on some open set containing $V(t)$ for all sufficiently small $t>0$. Then
    $$ \Div \bm F(\bm a) = \lim_{t\to 0^+} \frac{1}{|V(t)|} \iint_{S(t)} \bm F \cdot \dd \bm S,$$
    where $|V(t)|$ is the volume of $V(t)$. 
\end{corollaryenv}
\begin{proofenv}
    Write $\varphi = \Div \bm F$. Let $\varepsilon > 0$. Then it suffices to show that there exists $\delta > 0$ such that for any $t\in (0,\delta)$, we have
    $$ \left| \frac{1}{|V(t)|} \iint_{S(t)} \bm F \cdot \dd \bm S - \varphi(\bm a) \right| < \varepsilon. $$
    Since $\varphi$ is continuous at $\bm a$, there exists $\delta > 0$ such that
    $$ |\varphi(\bm x) - \varphi(\bm a)| < \frac{\varepsilon}{2},\quad \text{whenever } \| \bm x - \bm a \| < \delta. $$
    Then, for any $t\in (0,\delta)$, since $\varphi(\bm a)$ is constant, we have
    $$
    \varphi(\bm a) |V(t)| = \varphi(\bm a) \iiint_{V(t)} 1 \dd V = \iiint_{V(t)} \varphi(\bm a) \dd V.
    $$
    By Gauss's theorem we have
    $$ \begin{aligned}
        \iiint_{V(t)} \varphi(\bm a) \dd V
        =\iiint_{V(t)} \varphi(\bm x) + \iiint_{V(t)} \varphi(\bm a) - \varphi(\bm x) \dd V = \iint_{S(t)} \bm F \cdot \dd \bm S + \iiint_{V(t)} \varphi(\bm a) - \varphi(\bm x) \dd V
    \end{aligned} $$
    and thus
    $$ \left| \varphi(\bm a) |V(t)| - \iint_{S(t)} \bm F \cdot \dd \bm S \right| = \left| \iiint_{V(t)} \varphi(\bm a) - \varphi(\bm x) \dd V \right| \le \iiint_{V(t)} |\varphi(\bm a) - \varphi(\bm x)| \dd V \le |V(t)| \cdot \frac{\varepsilon}{2} < \varepsilon |V(t)|. $$
    Dividing both sides by $|V(t)|$, we have
    $$ \left| \frac{1}{|V(t)|} \iint_{S(t)} \bm F \cdot \dd \bm S - \varphi(\bm a) \right| < \varepsilon. $$
    This completes the proof.
\end{proofenv}
\begin{remarkenv}
    The corollary above gives us a coordinate free formula for the divergence of a vector field at a point. It also suggests the geometric interpretation of divergence: the divergence of a vector field at a point is the infinitesimal net outward flux per unit volume at that point.
\end{remarkenv}